\subsection{Definitions and notations}

Let G be a topological group operating continuously on the left (GT, III, §2, No.~4) in a locally compact space X; for \(s \in G\) and \(x \in X\) , let sx be the transform of x by s. We denote by \(\boldsymbol{\gamma}_{\mathbf{X}}(s)\) , or \(\boldsymbol{\gamma}(s)\) , the homeomorphism of X onto X defined by

\[
\gamma (s) x = s x.\tag{1}
\]

We have

\[
\boldsymbol {\gamma} (s t) = \boldsymbol {\gamma} (s) \boldsymbol {\gamma} (t).\tag{2}
\]

If \(f\) is a function defined on \(X\) , \(\boldsymbol{\gamma}(s)f\) will be defined by transport of structure, that is, by the formula \((\boldsymbol{\gamma}(s)f)(\boldsymbol{\gamma}(s)x) = f(x)\) ; in other words,

\[
\big (\boldsymbol {\gamma} (s) f \big) (x) = f (s ^ {- 1} x).\tag{3}
\]

If \(\mu\) is a measure defined on X, \(\gamma(s)\mu\) will also be defined by transport of structure, which leads to

\[
\langle f, \boldsymbol {\gamma} (s) \mu \rangle = \langle \boldsymbol {\gamma} (s ^ {- 1}) f, \mu \rangle \quad \text {   for   } f \in \mathscr {K} (\mathrm{X}).\tag{4}
\]

In other words,

\[
\int_ {\mathrm{X}} f (x) d (\boldsymbol {\gamma} (s) \mu) (x) = \int_ {\mathrm{X}} f (s x) d \mu (x).\tag{5}
\]

If A is a \((\gamma(s)\mu)\) -integrable set, then \(s^{-1}A\) is \(\mu\) -integrable and

\[
\left(\boldsymbol {\gamma} (s) \mu\right) (\mathrm{A}) = \mu \left(s ^ {- 1} \mathrm{A}\right).\tag{6}
\]

The measure \(\boldsymbol{\gamma}(s)\mu\) may also be defined as the image of \(\mu\) under \(\boldsymbol{\gamma}(s)\) . Instead of writing \(d(\boldsymbol{\gamma}(s)\mu)(x)\) , it is sometimes useful to write \(d\mu(s^{-1}x)\) ; the formula (5) then takes the following form:

\[
\int_ {\mathrm{X}} f (x) d \mu (s ^ {- 1} x) = \int_ {\mathrm{X}} f (s x) d \mu (x);
\]

the expression on the right may then be deduced from that on the left by `replacing \(x\) by \(sx\) .'

DEFINITION 1. --- Let \(\mu\) be a measure on X.

\begin{enumerate}
\def\labelenumi{\alph{enumi})}
\item
  \(\mu\) is said to be invariant under G if \(\pmb{\gamma}(s)\mu = \mu\) for every \(s\in \mathbf{G}\) .
\item
  \(\mu\) is said to be relatively invariant under \(\mathbf{G}\) if \(\pmb{\gamma}(s)\pmb{\mu}\) is proportional to \(\pmb{\mu}\) for every \(s\in \mathbf{G}\) .
\item
  \(\mu\) is said to be quasi-invariant under \(\mathbf{G}\) if \(\pmb{\gamma}(s)\pmb{\mu}\) is equivalent to \(\pmb{\mu}\) for every \(s\in \mathbf{G}\) .
\end{enumerate}

Remarks. --- 1) Assume \(\mu\) invariant. Then \(|\mu|\) , \(\mathcal{R}(\mu)\) , \(\mathcal{I}(\mu)\) are invariant. If \(\mu\) is real, then \(\mu^{+}\) and \(\mu^{-}\) are invariant.

\begin{enumerate}
\def\labelenumi{\arabic{enumi})}
\setcounter{enumi}{1}
\tightlist
\item
  Assume \(\mu\) relatively invariant and nonzero. There exists, for every \(s \in G\) , a unique complex number \(\chi(s)\) such that
\end{enumerate}

\[
\boldsymbol {\gamma} (s) \mu = \chi (s) ^ {- 1} \mu ,\tag{7}
\]

and the function \(\chi\) on \(\mathbf{G}\) is a representation of \(\mathbf{G}\) in \(\mathbf{C}^*\) , called the multiplier of \(\mu\) . The formula (5) then gives

\[
\int_ {\mathrm{X}} f (s x) d \mu (x) = \chi (s) ^ {- 1} \int_ {\mathrm{X}} f (x) d \mu (x),\tag{8}
\]

and formula (6) gives

\[
\mu (s \mathrm{A}) = \chi (s) \mu (\mathrm{A}).\tag{9}
\]

With the conventions made above, (7) may also be written

\[
d \mu (s x) = \chi (s) d \mu (x).\tag{10}
\]

\begin{enumerate}
\def\labelenumi{\arabic{enumi})}
\setcounter{enumi}{2}
\tightlist
\item
  Since \(|\pmb{\gamma}(s)\pmb{\mu}| = \pmb{\gamma}(s)(|\pmb{\mu}|)\) , to say that \(\pmb{\mu}\) is quasi-invariant amounts to saying that \(|\pmb{\mu}|\) is quasi-invariant.
\end{enumerate}

If \(\mu\) is quasi-invariant and \(\mu'\) is another measure on X equivalent to \(\mu\) , then \(\gamma(s)\mu'\) is equivalent to \(\gamma(s)\mu\) , hence to \(\mu\) , hence to \(\mu'\) , and so \(\mu'\) is quasi-invariant. To say that \(\mu\) is quasi-invariant under G therefore means that the class of \(\mu\) is invariant under G.

For \(\mu\) to be quasi-invariant, it is necessary and sufficient that the set of locally \(\mu\) -negligible subsets of X be invariant under G (Ch. V, §5, No.~5, Th. 2), or again that, for every \(\mu\) -negligible compact subset K of X and for every \(s \in \mathbf{G}\) , sK be \(\mu\) -negligible (loc. cit., Remark).

If \(\mu\) is quasi-invariant, then the support of \(\mu\) is invariant under G. In particular if G is transitive in X (A, I, §5, No.~5, Def. 6), this support is either empty (if \(\mu = 0\) ) or equal to X (if \(\mu \neq 0\) ).

Lemma 1. --- Let X, Y, Z be three topological spaces, with Y locally compact. Let \((x,y)\mapsto xy\) be a continuous mapping of \(X\times Y\) into Z, which defines a mapping \(x\mapsto u_{x}\) of X into \(\mathcal{F}(Y;Z)\) by the relation \(u_{x}(y)=xy\) . Let f be a continuous function on Z with values in R or in a Banach space, S the support of f, and \(\mu\) a measure on Y. Assume that for every \(x_{0}\in X\) , there exists a neighborhood V of \(x_{0}\) in X such that \(\bigcup_{x\in V}u_{x}^{-1}(S)\) is relatively compact in Y. Then:

\begin{enumerate}
\def\labelenumi{\alph{enumi})}
\tightlist
\item
  for every \(x \in X\) , \(f \circ u_{x}\) is continuous on Y, with compact support;
\item
  the mapping \(x \mapsto \int_{Y} f(xy) d\mu(y)\) , which is defined by a), is continuous on X.
\end{enumerate}

The assertion a) is obvious. Let us prove b). Since continuity is a local property, we may reduce to the case that \(\bigcup_{x\in X}u_x^{-1}(S)\) is contained in a compact subset \(Y'\) of Y. Since the function \((x,y)\mapsto f(xy)\) is continuous on \(X\times Y\) , \(f\circ u_x\) tends to \(f\circ u_{x_0}\) uniformly on \(Y'\) as \(x\) tends to \(x_0\) (GT, X, §3, No.~4, Th. 3), therefore \(\mu(f \circ u_x)\) tends to \(\mu(f \circ u_{x_0})\) . Whence the lemma.

Let us now return to the previous notations.

PROPOSITION 1. --- Assume that G is locally compact. Let \(\mu\) be a nonzero relatively invariant measure on X. Then its multiplier \(\chi\) is a continuous function on G.

For, let \(f \in \mathcal{K}(\mathbf{X})\) , \(S\) the support of \(f\) , \(s_0\) a point of \(G\) , and \(V\) a compact neighborhood of \(s_0\) in \(G\) ; then, the set

\[
\bigcup_ {s \in V} \gamma (s) ^ {- 1} (S) = V ^ {- 1} S
\]

is compact in X; by Lemma 1 and formula (8), \(\chi(s^{-1})\langle\mu, f\rangle\) depends continuously on \(s\) ; if \(f\) is chosen so that \(\langle\mu, f\rangle \neq 0\) , one sees that \(\chi\) is continuous.

Now let G be a topological group operating continuously on the right in a locally compact space X; for \(s \in G\) and \(x \in X\) , let xs be the transform of x by s. We denote by \(\delta_{\mathbf{X}}(s)\) , or \(\delta(s)\) , the homeomorphism of X defined by

\[
\delta (s) x = x s ^ {- 1}.\tag{1'}
\]

We have

\[
\boldsymbol {\delta} (s t) = \boldsymbol {\delta} (s) \boldsymbol {\delta} (t).\tag{2'}
\]

By transport of structure, one defines the action of \(\boldsymbol{\delta}(s)\) on functions and measures on X:



\[
\big (\boldsymbol {\delta} (s) f \big) (x) = f (x s)\tag{3'}
\]

\[
\langle f, \delta (s) \mu \rangle = \langle \delta (s ^ {- 1}) f, \mu \rangle\tag{4'}
\]

\[
\int_ {\mathrm{X}} f (x) d \bigl (\pmb {\delta} (s) \mu \bigr) (x) = \int_ {\mathrm{X}} f (x s ^ {- 1}) d \mu (x)\tag{5'}
\]

\[
\big (\delta (s) \mu \big) (\mathrm{A}) = \mu (\mathrm{A} s).\tag{6'}
\]

We agree to write \(d\mu(xs)\) in place of \(d(\pmb{\delta}(s)\mu)(x)\) , and \((5')\) then takes the form

\[
\int_ {\mathrm{X}} f (x) d \mu (x s) = \int_ {\mathrm{X}} f (x s ^ {- 1}) d \mu (x).
\]

One defines in an analogous manner the measures on X invariant, relatively invariant and quasi-invariant under G. If \(\mu\) is relatively invariant, its multiplier \(\chi\) is defined by the formulas



\[
\boldsymbol {\delta} (s) \mu = \chi (s) \mu\tag{7'}
\]

\[
\int_ {\mathrm{X}} f (x s) d \mu (x) = \chi (s) ^ {- 1} \int_ {\mathrm{X}} f (x) d \mu (x)\tag{8'}
\]

\[
\mu (\mathrm{A} s) = \chi (s) \mu (\mathrm{A})\tag{9'}
\]

\[
d \mu (x s) = \chi (s) d \mu (x).\tag{10'}
\]

If one regards the group \(G^{0}\) opposite to G as operating in X by \((x, s) \mapsto xs\) , then \(\mu\) is relatively invariant under \(G^{0}\) with the same multiplier \(\chi\) .

Finally, let \(\mathbf{G}\) be a locally compact group. It operates on itself by left and right translations, according to the formulas \(\pmb{\gamma}(s)x = sx\) , \(\pmb{\delta}(s)x = xs^{-1}\) . Then

\[
\boldsymbol {\gamma} (s) \boldsymbol {\delta} (t) = \boldsymbol {\delta} (t) \boldsymbol {\gamma} (s).\tag{11}
\]

All of the foregoing is applicable here, thus we have on G the concepts of measures that are left-invariant, right-invariant, relatively left-invariant, relatively right-invariant, left quasi-invariant, right quasi-invariant (however, see Nos. 8 and 9).

The mapping \(x \mapsto x^{-1}\) is a homeomorphism of \(G\) onto \(G\) . For every function \(f\) on \(G\) , we define the function \(\overset{\vee}{f}\) on \(G\) by

\[
\stackrel {\vee} {f} (x) = f (x ^ {- 1}).\tag{12}
\]

For every measure \(\mu\) on \(\mathbf{G}\) , we define the measure \(\breve{\mu}\) by

\[
\stackrel {\vee} {\mu} (f) = \mu (\stackrel {\vee} {f}) \quad \text { for } f \in \mathcal {K} (\mathrm{G}).\tag{13}
\]

In other words,

\[
\int_ {\mathrm{G}} f (x) d \mu^ {\vee} (x) = \int_ {\mathrm{G}} f (x ^ {- 1}) d \mu (x).\tag{14}
\]

If \(\mathbf{A}\) is a \(\check{\mu}\) -integrable set, then \(\mathbf{A}^{-1}\) is \(\mu\) -integrable and

\[
\check {\mu} (\mathrm{A}) = \mu (\mathrm{A} ^ {- 1}).\tag{15}
\]

We agree to write \(d\mu(x^{-1})\) in place of \(d\check{\mu}(x)\) , and (14) then takes the form

\[
\int_ {\mathrm{G}} f (x) d \mu (x ^ {- 1}) = \int_ {\mathrm{G}} f (x ^ {- 1}) d \mu (x).
\]

